% Discussion
\section{Discussion}

Our calibrated transcriptomic classifier establishes a scalable framework for inferring porcine developmental stages across eight tissues. The approach harmonises chronological ages into biologically meaningful stages and delivers high predictive performance with tight confidence intervals. The methodology allows laboratories to stage archived or prospective samples without relying on fragile morphological features, providing practical value to both agricultural management and translational research.

The enrichment and marker analyses demonstrate that the models capture bona fide developmental programs rather than technical artefacts. Tissue-resolved signatures recover known regulators of adipogenesis, neurodevelopment, hepatic maturation, and myogenesis, while also flagging respiratory and intestinal pathways that align with phenotypic maturation timelines. These results provide a molecular atlas that can inform targeted interventions, such as timing nutritional supplementation or assessing readiness for organ-harvesting protocols.

A particularly significant finding is our validation of conserved biomarkers between pig and human developmental programs. The observation that canonical human developmental biomarkers show expression patterns consistent with published human developmental studies provides compelling evidence for the utility of swine as a model organism for human development and disease. This qualitative concordance, based on literature-reported functions, is not limited to a single biological system but spans diverse developmental processes: hepatocyte differentiation (ALB), CNS myelination (MBP), erythroid maturation (HBB), adipogenesis (PPARG), myogenesis regulation (MSTN), alveolar development (SFTPC), intestinal homeostasis (LGR5), and gonadal development (DMRT1). This extensive cross-species conservation strongly suggests that developmental timing mechanisms identified in our porcine model may directly inform human pediatric medicine, regenerative biology, and age-related disease research. The ability to accurately stage developmental processes using molecular markers rather than morphological features enables more precise timing of interventions and reduces variability in translational studies, making the pig an even more valuable model for biomedical research.

Future research directions include extending the classifier to prenatal stages, integrating single-cell references for enhanced resolution, and benchmarking against independent cohorts to further strengthen translational readiness. The tissue-specific nature of developmental programs, as revealed by our analyses, underscores the importance of maintaining tissue-resolved models rather than attempting to force generalisation across fundamentally distinct developmental trajectories.

In conclusion, this molecular staging system represents a significant advance in porcine developmental biology. By providing objective, quantitative staging independent of morphological assessment, the framework enables more precise experimental designs, reduces inter-observer variability, and facilitates retrospective analysis of archived samples. As transcriptomic resources continue to expand, this approach can be refined to achieve even finer temporal resolution, ultimately supporting both fundamental research and practical applications in agriculture and biomedicine.